% !TEX root = ../linnik399.tex
\section{The computations and their verification}\label{sec:verification}

The finite computation has three logical components. The zero-location checks justify
the case analysis; enclosures of analytic quantities supply the leaf coefficients;
and integer certificates bound the resulting programs. The exterior ranges require
separate computations. A check of one component does not by itself verify the others.
This section records the evidence for each and the scope of the Lean formalization.

Computer assistance is well established in explicit analytic number theory, for instance in the
ternary Goldbach problem \cite{Helfgott2015ternary} and in verifications of the Riemann hypothesis
\cite{Platt2016numerical,PlattTrudgian2021riemann}, and linear programs with rigorously checked
bounds are central to the proof of the Kepler conjecture \cite{ObuaNipkow2009flyspeck,
SolovyevHales2011efficient,Hales2017formal}. In our computations, the real coefficients and
margins are enclosed by interval arithmetic \cite{Moore1966interval,MooreKearfottCloud2009introduction,
Rump2010verification,Tucker2011validated}. It is implemented with the interval context of mpmath
\cite{mpmath2023library}, at a working precision of at least $50$ digits, and with a vectorized
binary64 interval module that widens the result of every operation outward by one unit in the
last place and evaluates $\exp$, $\sin$ and $\cos$ by Taylor polynomials with explicit remainder
bounds, not by the platform's library; the Lean checkers of \cref{subsec:lean} use their own
verified rational interval arithmetic. Floating-point linear programming (with the HiGHS solver
\cite{HuangfuHall2018parallelizing} in SciPy \cite{Virtanen2020scipy}) is used only to propose dual
multipliers, and every certificate is checked in exact integer arithmetic. Independent
numerical audits use multiple-precision arithmetic \cite{mpmath2023library}.
The supplementary floating-point checks of certain published table entries are
identified separately in \cref{subsec:locconventions}. \Cref{tab:status} summarizes, for each
component of the proof, whether it is published or new, how it was obtained and how it was
checked.

\begin{table}[tp]
\centering
\footnotesize
\setlength{\tabcolsep}{4pt}
\begin{tabular}{@{}>{\raggedright\arraybackslash}p{\dimexpr0.25\linewidth-6.4pt\relax}
                >{\raggedright\arraybackslash}p{\dimexpr0.15\linewidth-6.4pt\relax}
                >{\raggedright\arraybackslash}p{\dimexpr0.11\linewidth-6.4pt\relax}
                >{\raggedright\arraybackslash}p{\dimexpr0.15\linewidth-6.4pt\relax}
                >{\raggedright\arraybackslash}p{\dimexpr0.34\linewidth-6.4pt\relax}@{}}
\toprule
Component & Source & Kind & Method & Checked by\\
\midrule
Positivity criterion and detection (\cref{inp:criterion,prop:detect}) &
  Published \cite{Xylouris2011nullstellen} & Analytic & -- &
  Lean, with the criterion as a hypothesis\\
Envelope and first-family bounds (\cref{sec:costs}) &
  New, after \cite{Xylouris2011nullstellen} & Analytic & Interval arithmetic &
  Hand proof; a hypothesis in Lean, which checks only the enclosures\\
Far density (\cref{inp:far}) &
  Published \cite{Xylouris2011nullstellen} & Analytic & -- &
  Lean hypothesis; profile conditions proved in Lean\\
Far constants $V$, $w$ (\cref{sec:far}) &
  New profiles & Numerical & Interval arithmetic &
  Multiple-precision audit; Lean numeric checker\\
Printed zero-location tables (\cref{prop:parentrows}) &
  Published \cite{HeathBrown1992zero,Xylouris2011nullstellen} & Analytic & As printed &
  Compared with the pages; Heath-Brown's Tables~4 and~7 recomputed in floating point\\
New location rows (\cref{sec:location}) &
  New & Both & Interval arithmetic &
  Replay; the $94$ complex rows by a separate program\\
Conditions \cite[(4.29), (4.34)]{Xylouris2011nullstellen} &
  Published; new check & Numerical & Interval arithmetic &
  Separate programs\\
Conductor refinement $\frac1{108}$ (\cref{subsec:conductor}) &
  New & Analytic & -- & Hand proof\\
Graded near lemma, response lemma, threshold form (\cref{sec:near}) &
  New & Analytic & -- & Lean (Comparator, two kernels)\\
Family, shifted and graded rows (\cref{sec:rows}) &
  New & Numerical & Interval arithmetic &
  Lean row checker ($16{,}842$ rows); multiple-precision audit\\
Two-test rows (\cref{subsec:twotest}) &
  New, after \cite{Xylouris2011nullstellen} & Both & Interval arithmetic &
  Hand proof; replay; exact comparison with a separate implementation on every certificate
  interval (constants: that implementation's interval code only)\\
Certificate soundness (\cref{thm:certificate}) &
  New & Analytic & -- & Lean (Comparator, two kernels)\\
Leaf certificates (\cref{def:certificate}) &
  New & Numerical & Exact integer &
  Replay; independent checker; Lean checker\\
Leaf data (\cref{subsec:leafmodel}) &
  New & Numerical & Interval arithmetic &
  Multiple-precision audits; Lean numeric checker (count-type integers: audit only)\\
Case tree and realization (\cref{sec:casetree}) &
  New & Analytic & -- & Hand proof; replay\\
Tiny exceptional zero (\cref{prop:tiny}); empty rectangle &
  Published \cite{HeathBrown1990siegel}; $W=0$ & Analytic & -- & As printed\\
Small exceptional zero (\cref{thm:small}) &
  New, from lemmas of \cite{HeathBrown1992zero} & Both & Interval arithmetic &
  Hand proof; exterior replay; margins recomputed independently\\
Large first zero (\cref{thm:large}) &
  New & Numerical & Exact integer &
  Replay's exact checker; independent checker; Lean certificate (native and kernel), numeric and
  row checkers\\
Assembly (\cref{sec:join}) &
  New & Analytic & -- & Hand proof\\
\bottomrule
\end{tabular}

\medskip
\begin{minipage}{\linewidth}
\footnotesize
The interval arithmetic is described at the beginning of this section. Floating point is used
only to propose parameters, dual multipliers and sieve heights, which need no justification, and
to recompute published tables.
\end{minipage}
\caption{Status of the components. Published inputs are used as printed, and those used in the
Lean theorems enter them as hypotheses; ``Both'' means an analytic argument with numerical
margins; ``hand proof'' means a conventional proof in this paper. The checks are described in
\cref{sec:verification}.}
\label{tab:status}
\end{table}

\subsection{The data}
The certificates are stored in three compressed JSON record files. The inside and
outside files are indexed by root case; the supplementary file records the $233$
supplementary leaves of $199$ inside roots (\cref{subsec:refinement}). Their sizes and SHA-256 hashes are:
\begin{center}
\small
\begin{tabular}{lr}
\toprule
File & Bytes \\
\midrule
\texttt{inside.jsonl.gz} & $298{,}536{,}973$ \\
\multicolumn{2}{l}{\scriptsize\texttt{539362a03e7b330dad03b33d6b5037871cfcbbf7b0973e7976fccf12509fe63c}}\\[2pt]
\texttt{nodes.jsonl.gz} & $159{,}624{,}716$ \\
\multicolumn{2}{l}{\scriptsize\texttt{b715ac735b6976da17a2671667ecfdf07c3f19b8ab74ae5f0533dbcf6b2a6cfc}}\\[2pt]
\texttt{outside.jsonl.gz} & $8{,}626{,}435$ \\
\multicolumn{2}{l}{\scriptsize\texttt{a53d4d5793f71c92647a7f22022110bdc8135cbda1f2575d1dc0c1f487f5ac1d}}\\
\bottomrule
\end{tabular}
\end{center}
The first holds the $2768$ inside roots with their $2913$ other leaves, the second the $233$
supplementary leaves, and the third the $1685$
outside roots with their $1642$ leaves. A leaf record contains the parameters of its near rows
and its certificate tree; it does \emph{not} contain the leaf program itself, which the verifier
regenerates from the case tree. Some leaves are split into subcases (\cref{subsec:subdivision}), each with
its own certificate tree, so the $4788$ leaves carry $5591$ certificate trees in all: $3949$
inside (including the supplementary leaves) and $1642$ outside.

\subsection{Exact verification}
All numbers entering a leaf program are enclosed with outward-rounded interval arithmetic and
stored as integers scaled by $S=10^{16}$, rounded in the safe direction (\cref{sec:lp}). The
verification of a certificate uses only integer arithmetic. The following checks have been run.
\begin{enumerate}[(1)]
  \item \emph{The replay.} A verifier traverses the case tree from its roots. It re-checks every
    split, location node and exclusion (regenerating the location rows, except that the table
    of complex polynomial rows and the side conditions \cite[(4.29), (4.34)]{Xylouris2011nullstellen} are
    checked by separate programs).
    At every leaf it regenerates the leaf program from the specification and the tree,
    independently of the stored certificate, and checks the certificate exactly. It passes on all
    $4453$ roots: $4788$ leaves, $4{,}196{,}879$ boxes, and largest certified value
    $0.9999998223\ldots<1$ (inside) and $0.9999942478\ldots$ (outside).
  \item \emph{An independent checker.} A second verifier was written from \cref{sec:lp} alone,
    without reading or using the checking code of the first, and it uses only exact integer and
    rational arithmetic. It takes the leaf data from the same generator as the replay (the leaf
    data are checked separately, in (4) and in \cref{subsec:lean}). It accepts every certificate of
    the corpus ($29{,}397{,}336$ relaxation cases in all) and that of \cref{thm:large}, and for every
    certificate tree its value and number of boxes are exactly those of the first verifier. In a
    negative test on ten leaves, including the tightest, it rejected all $308$ certificates made
    invalid by deliberate forgeries, among them duals lowered, and coefficients raised, by the
    smallest amount that breaks \eqref{eq:dualfeasible}; it accepted all $82$ controls perturbed by
    one unit less. It replaces an earlier independent checker, whose code was not preserved and
    which had also accepted every certificate of the corpus.
  \item \emph{A mutation test.} On the $19$ leaves of three of the tightest inside roots, $133$
    perturbations of the regenerated leaf data (far budget, objective, first-family charge,
    correlation term, features), of certificate multipliers and of the allowance $5\eta$, each in the
    unfavourable direction, were all rejected by the replay.
  \item \emph{Audits of the constants.} For $38$ leaves of all kinds, the $19{,}616$ objective and
    $19{,}616$ far coefficients, the tail, hidden and second-family columns, $F$, $\mathrm{first}$ and
    all $120$ family, shifted and graded near rows were recomputed from their definitions in
    multiple-precision arithmetic by a program sharing no code with the builders; all lie on the
    safe side of their true values. (The program of this audit was not preserved; the numeric and
    row checkers of \cref{subsec:lean} now check the same quantities for every leaf, except the
    two-test rows.) Separately, for $117$ roots ($100$ of them chosen at random),
    all $530{,}942$ column and budget integers of their $154$ leaves (objective, far, count,
    hidden-count and second-family coefficients, $F$, $\mathrm{first}$ and $\mathrm{final}$) were
    recomputed from their defining formulas in $50$-digit arithmetic and agree exactly with the
    stored ones.
  \item \emph{The two-test rows.} A separate program replays the $37$ roots whose leaves use the
    row of \cref{subsec:twotest} ($31$ with a single test, $2$ with a mixture, $4$ paired). It
    compares the row as the certificates use it with a separate implementation of the same
    inequality, with its own builder and its own certificate checker. The
    comparison runs on every threshold interval that the $37$ certificates use ($1079$ intervals,
    from $23{,}341$ boxes) and on $11{,}211$ seeded random intervals, in the first-order relaxation
    and in both tangent relaxations (for these, the tangent rule applied to the exact inequality of
    that implementation). The exact costs and budgets are those of that inequality divided by $D$,
    and the integers used are rounded in the safe direction. The row's integers are regenerated by
    the interval routines of that implementation and agree with the leaves; every one-unit
    perturbation of them is detected. These constants rest on those interval routines alone.
\end{enumerate}

\emph{The exterior regimes.} The margins of \cref{thm:small} were regenerated from their
definitions and checked in interval arithmetic by a separate replay of the exterior components,
and they have been recomputed independently in multiple-precision interval arithmetic. The
program of \cref{thm:large} is a leaf program in the sense of \cref{def:leafprogram}. Its
objective and far integers are regenerated and compared with the stored ones, and its near row is
built with the normalization of \cref{subsec:rowdata}. The exact checking code of the replay (1)
and the independent checker (2) accept its certificate, and so do the verified checkers of
\cref{subsec:lean}, which also accept its numeric data and its near row.

\subsection{Formal verification}\label{subsec:lean}
Several parts of the argument have been formalized in the Lean~4 proof assistant
\cite{deMouraUllrich2021lean} with the Mathlib library \cite{mathlib2020lean}; see
\cite{Avigad2024mathematics} for the role of such formalizations. Earlier formalizations in
analytic number theory include the prime number theorem
\cite{AvigadDonnellyGrayRaff2007formally,Harrison2009formalizing,KontorovichTao2024pnt} and
Dirichlet's theorem \cite{Stoll2024dirichlet}. The formalization uses Lean \texttt{v4.35.0-rc3} with
the corresponding Mathlib release, and contains about $24{,}000$ lines and about $800$ theorems and
lemmas. The $18$ statements of the near lemma and the certificate checker were checked against
their proofs with the Comparator tool \cite{LeanFRO2025comparator}, which replays the proofs in the
Lean kernel and in the independent kernel nanoda \cite{Bailey2020nanoda}; the leaf-level and near-row
theorems were checked by the Lean kernel. All proofs are complete and use only the
three standard axioms (propositional extensionality, the axiom of choice, and the soundness of
quotients). The published inputs of \cref{sec:inputs} enter as explicit hypotheses, each stated
no more strongly than its printed source. The following results are formally verified.
\begin{enumerate}[(a)]
  \item The graded near-density lemma (\cref{thm:graded}), from the local explicit formulas,
    Burgess's estimate and Graham's estimate as hypotheses. This includes the response lemma,
    the threshold form, and the entry bounds of \cref{sec:rows}.
  \item The soundness of the certificate checker of \cref{sec:lp}: an accepted certificate
    proves that every feasible point of the leaf program has value below $1$.
  \item The passage from a certified leaf with valid data to primes. From Xylouris's criterion
    and his far-density lemma, stated as printed, and the localized envelope of
    \cref{prop:envelope}, stated as a hypothesis, a leaf whose certificate and data pass the checks
    yields a prime $q^{A}<p<q^{L}$ in every reduced class, for every modulus whose zeros satisfy
    the leaf's zero-level hypotheses. The count, hidden-count and second-family integers and
    $\mathrm{final}$ enter as they are (they were recomputed in the audit of (4) above).
  \item Interval arithmetic with proved inclusion, including the exponential function, and
    enclosures of the functions entering the objective and far data of the leaves. A numeric
    checker is proved to imply that each leaf's objective and far coefficients, far budget and
    first-family charge are valid.
  \item Heath-Brown's Conditions~1 and~2 for the parabolic test functions of \cref{sec:rows}.
  \item The near rows: each family, shifted and graded near row is a concrete instance of
    \cref{thm:graded}. A row checker is proved sound, including the three family terms that keep a second zero and the constant
    $C_Z$ of \cref{sec:rows}. A checked row is a valid row of the leaf program under
    zero-level hypotheses only, and this is composed with (c). (The formal version of
    \cref{cor:zeroform} assumes that the numerators $D(\delta_j)-d_1+e_j$ themselves are positive,
    which is stronger than what the proof in \cref{sec:near} needs; the row checker verifies this
    for every checked entry.)
\end{enumerate}
The certificate checker has a Mathlib-free copy that is proved to compute the same value; the
numeric checker and the row checker are Mathlib-free themselves and are proved sound directly.
Run as compiled programs, the certificate and numeric checkers accept every certificate tree and
the objective, far and first-family data of every leaf in the corpus: $5591$ trees and
$4{,}196{,}879$ boxes.
The compiled row checker accepts every near row of the corpus except the $37$ two-test rows of
\cref{subsec:twotest}, which it does not cover: $16{,}842$ rows with $11{,}135{,}323$ entries. The same
three checkers accept the leaf program of \cref{thm:large}: its certificate ($36$ boxes), its
objective and far data, and its near row ($150$ entries). So (c) and (f) apply to it as to the leaves of the
case tree. In addition, seven certificates are checked by evaluation in the Lean kernel itself:
six of the case tree ($122$ boxes) and that of \cref{thm:large} ($36$ boxes).

The published inputs (\cref{sec:inputs,sec:location,sec:exteriors}) enter the formal
statements as hypotheses. The following parts have conventional proofs in this paper
but are \emph{not} formally verified. Their numerical checks are as described here and
in the relevant sections; in particular, a multiple-precision audit is distinct from
a formally verified enclosure.
\begin{itemize}
  \item the per-character costs of \cref{sec:costs}: the localized envelope, which enters as a
    hypothesis, and the first-family bounds (the enclosures of $J_{\mathrm{old}}$ and $J_{\mathrm{new}}$
    are verified, but not the inequalities they bound);
  \item the conductor coefficient $\frac14$ for real characters and the refinement of
    \cref{subsec:conductor};
  \item the zero-location arguments and the case tree (\cref{sec:location,sec:casetree}), that is, the
    statement that every configuration of zeros falls into a leaf whose zero-level hypotheses it
    satisfies, including the count, hidden-count and second-family constraints;
  \item the exterior regimes (\cref{sec:exteriors}): the range of \cref{prop:tiny,thm:small}, and,
    for \cref{thm:large}, the statement that every configuration with $\lambda_1\ge1.5$ satisfies the
    zero-level hypotheses of its program;
  \item the two-test row of \cref{subsec:twotest}, which is checked by the replay and compared
    exactly with a separate implementation in (5), and whose constants were computed only by the
    interval routines of that implementation;
  \item the count-type integers of the leaf programs (checked in the audit of (4));
  \item the assembly of \cref{sec:join}.
\end{itemize}
A compiled program is trusted to compute what it is proved to compute; its correctness then rests
on the Lean compiler and runtime and on the unverified routines that read the data files. The Lean
proofs were written by AI agents, like the rest of this work (see the note after the abstract).
The Lean kernel checks the proofs, but whether each formal statement says what the corresponding
statement of this paper says has so far been reviewed only by AI agents.

\subsection{Data availability}
The development repository contains the certificates, the programs that generate and
check them, the verification reports, and the Lean formalization:
\begin{center}\url{https://github.com/enaslund/linniks-constant}.\end{center}
At the date of this revision the repository is private. A public, versioned archive
of these materials is needed before submission so that readers can reproduce the
computations and identify the precise data used in the proof. This paper and the Lean
formalization are publicly available in the repository prepared for the Palomar registry of
formalized mathematics:
\begin{center}\url{https://github.com/enaslund/linniks-constant-3.99}.\end{center}

The repository includes instructions for rerunning the replay, mutation tests,
numerical audits, and formal checks. On the machine used,
the replay took $1.7$ hours of wall-clock time with parallel workers, and the compiled Lean checks
of all certificates, leaf data and near rows took about $5$ hours of wall-clock time. One caveat
concerns exact reproducibility: the sieve heights of the graded rows (\cref{subsec:rowdata}) are
regenerated in binary64 floating point and are not stored in the records. Any nonnegative heights
are admissible, so this does not affect soundness, but on a platform whose elementary functions
round differently the regenerated rows, and hence the certificates that use them, could differ and
would have to be recomputed. Thus the reproducibility requirement concerns the
particular stored certificates; admissibility of the analytic sieve weight alone does
not guarantee that a different regenerated row will pass those certificates.
